\title{\Large \bf Analyzing  the transference of the coalgebra
structure on the homology of CDGAs}
\author{\normalsize M.J. Jim\'{e}nez, P. Real and B. Silva\\
\normalsize Dpto. Matem\'{a}tica Aplicada I (Universidad de Sevilla)\\
\normalsize Avda. Reina Mercedes, s/n. CP: 41012 Sevilla (Spain)\\
\normalsize  e-mail: majiro@euler.fie.us.es, real@cica.es,
silva@cica.es}
\documentstyle[emlines2,bezier,11pt]{article}
\textheight=615pt
\textwidth=425pt
\oddsidemargin=0pt
\headheight=0pt
\headsep=7pt
\hoffset=0pt
\voffset=0pt
\topmargin=10pt
\marginparsep=10pt
\marginparwidth=30pt
\footheight=10pt
\footskip=40pt
\parskip=.5cm
\setlength{\baselineskip}{8ex}
\newcommand{\ra}{\rightarrow}
\newcommand{\la}{\leftarrow}
\newcommand{\scst}{\scriptscriptstyle}
\newcommand{\ot}{\otimes}
\newcommand{\Z}{\Bbb Z}
\newcommand{\pa}{\partial}
\newcommand{\al}{\alpha}
\newcommand{\be}{\beta}
\newtheorem{prp}{Proposition}[section]
\newtheorem{lmm}[prp]{Lemma}
\newtheorem{thr}[prp]{Theorem}
\newtheorem{dfn}[prp]{Definition}
\newtheorem{rmr}[prp]{Remark}
\newtheorem{crl}[prp]{Corollary}
\newtheorem{ejemplo}[prp]{Example}
\newtheorem{notacion}[prp]{Notation}
\hyphenation{pro-duc-to}
\begin{document}
\bibliographystyle{alpha}
\sloppy
\maketitle
\vspace{0.5cm}

\begin{abstract}

    Our motivation here is to analyze the $A_{\infty}$-coalgebra structure
 $(\Delta_2, \Delta_3, \ldots)$
of the small homological model $H$ obtained in \cite{Moscu}, for any commutative differential
graded algebra (or briefly CDGA) $A$. More precisely, making use
of the facts that $H$ is a CDGA and
that the morphism $\Delta_2: H \rightarrow
H\ot H$ satisfies a compatibility condition with regard to the product
on $H$, we design an  algorithm testing
the coassociativity of
 $\Delta_2$. Considerations concerning the complexity of this
 algorithm are given.

\end{abstract}

\section{Extended Abstract}



The description of efficient algorithms for homological computation
can be considered as a very important question in Homological
Algebra, in order to use those processes mainly in the resolution
of problems on Algebraic Topology. We are interested in
obtaining the $A_{\infty}$-coalgebra structures of small
homological models of Eilenberg-MacLane spaces. These spaces are a
kind of prime spaces in homotopy theory and its homology and
cohomology are essential in order to compute homotopy invariants
of spaces and cohomology operations. In order to progress in this
way, a first step is the analysis of the transference of the
coalgebra structure in the homology of CDGAs.

  Our main technique is the use of Homological Perturbation
  Theory which
 is often used to replace given chain
complexes by homotopic, smaller and more readily computable
chain complexes (\cite{Gug72}, \cite{GL89},\cite{GLS91}). The Basic
Perturbation Lemma (BPL) states
that given a contraction
\begin{equation}
(f,g,\phi):(N,d_{\scst N})\Rightarrow (M,d_{\scst M})
\label{contraction} \end{equation} of chain complexes (with
$fg=1_{\scst M}$ and $1_{\scst N}-gf = d_{\scst N}\phi + \phi
d_{\scst N}$) and a perturbation $\delta$ of $d_{\scst N}$ (that
is, $(d_{\scst N}+\delta)^{2}=0$), under suitable conditions there
exists a perturbation $d_{\delta}$ of $d_{\scst M}$ such that
$H_{*}(M,d_{\scst M}+d_{\delta})=H_* (N,d_{\scst N}+ \delta)$.


 Perturbation results about preservation of
structures (DG-algebras, DG-coalgebras, Lie algebras, \ldots) have
been largely considered (\cite{GLS91} and \cite{HK}). The technique for obtaining these
results is to determine under what
conditions the BPL allows the
preservation of the data structures.

 In the case of algebra contractions, the study of the preservation
 of the multiplicative structure of the first algebra via the
  contraction
leads to the definition of a special type of contractions:
{\em semi-full algebra contractions}.
The semi-fullness is an hereditary
property under composition, tensor product and perturbation of
contractions (which are new contractions).

Apart from the intrinsic interest of the transference
problem, when we deal with DG-algebras or DG-coalgebras,
these kind of results
allow us to determine and
construct essential algebra or coalgebra structures
on various chain complexes which have been constructed via
perturbation. For example, the
homology theory of commutative DGA-algebras is entirely examined in
terms of semifull algebra contractions in \cite{Moscu},
and  ``small''  $p$-local homological models of
 reduced bar constructions  of
twisted tensor product of Cartan's elementary complexes are
obtained in \cite{NY}, due to the fact that all the contractions
appearing there, are of this type. These examples will be our
starting point in this paper, being the transference  of the
coalgebra structure in  the previous
contractions our main objective.


  Let us begin by explaining the algorithm for computing small
  homological models of CDGAs.

It is commonly known that every commutative DGA  $A$ ``factors''
up to homotopy equivalence into a twisted
tensor product of exterior and polynomial
algebras endowed with a
differential-derivation (see, for example, \cite{MCL}); in
the sense that there exists an homomorphism
connecting both structures, which induces an isomorphism in
homology. This last TTP is called a {\em free model} for $A$.

 In fact, the object we start from is the free model of a
finite (as an algebra) CDGA
$A$.

The principal
goal that arises is to obtain a ``chain'' of contractions
starting at the
reduced bar construction $\bar{B}(A)$ and ending up at a
commutative
DGA-algebra that is free and of finite type as graded module.

Three {\it almost-full algebra contractions} (i.e., semifull algebra
contractions endowed with multiplicative projections) are  used in
order to find the structure of graded module of a homological model for the
DGA-algebra $A$:
\begin{itemize}
\item The contraction defined in \cite{EM54} from $\bar{B}(A\ot A')$ to
$\bar{B}(A)\ot \bar{B}(A')$, where $A$ and $A'$ are two commutative
augmented DGAs; which is briefly denoted by $C_{\scst \bar{B}\ot}$.

Though Eilenberg and Mac Lane only set explicit formulae for the projection
and inclusion morphisms, an explicit formula of the recursive definition of
the homotopy operator given by Eilenberg and Mac Lane in \cite{EM53} is
established in \cite{RUBIO} (see also \cite{Rea96a} for a proof).



Given a tensor product $\ot _{i\in I}A_i$ of augmented CDGA-algebras, a
contraction from $\bar{B}(\ot_{i\in I}A_i)$ to $\ot_{i\in I}\bar{B}(A_i)$ is
easily determined, by applying $C_{\scst \bar{B}\ot}$ several times in
a suitable way. This new contraction is also denoted by $C_{\scst
\bar{B}\ot}$.



\item The isomorphism of DGA (hence, a contraction) $C_{\scst
\bar{B}E}$ from $\bar{B}(E(u, 2n+1))$ to $\Gamma (\sigma (u), 2n+2)$, also
described in \cite{EM54}. Note that the generator $v$ of the previous
divided power algebra is denoted by $\sigma (u)$ where $\sigma$ is the Cartan's suspension
operator.

Last isomorphism might be considered as a {\it full algebra contraction}
(that is, an almost-full algebra contraction endowed with an algebra
homotopy operator), in an obvious way.

\item The contraction $C_{\scst \bar{B}P}$ from $\bar{B}(P(y,2n))$ to
$E(\sigma (y),2n)$ (see \cite{EM54}). The generator of the exterior algebra is
denoted by $\sigma (y)$ since both elements correspond  to each other
by the projection and inclusion of the contraction.
\end{itemize}

All the previous contractions have already been described for
decades, so we will not explain them further.


Thanks to these three contractions, it is possible to establish, by
composition, the following semifull algebra contraction $C_0 =(f,
g, \phi)$:
$$\bar{B}(\ot_{i\in I} A_i) \; \Rightarrow \; \ot_{i\in I} \bar{B}(A_i) \;
\Rightarrow \; \ot_{i\in I} HBA_i,$$
where $HBA_i$ represents an exterior or a divided polynomial algebra,
depending on the fact that $A_i$ is a polynomial or an exterior algebra.
Obviously, the product on $\ot_{i \in I} HBA_i$ is the natural one.


The next step is perturbing $C_0$, in order to obtain an homological model
of the initial TTP $\tilde{\ot}^{\rho}_{i\in I} A_i$. The morphism $\rho$
adds a perturbation-derivation $\delta$ onto the tensor
differential of $\bar{B}(\ot_{i \in I} A_i)$.

Therefore, a new semifull algebra
contraction $(f_{\delta},g_{\delta},\phi_{\delta})$ is obtained:
$$\bar{B}(\tilde{\ot}^{\rho}_{i\in
I} A_i) \stackrel{(C_0)_{\delta}} {\Rightarrow} (\ot_{i \in I}
HBA_i,d_{\delta}),$$
where the differential $d_{\delta}$ is determined by the perturbation
process. That means that $HBA=\ot_{i \in I} (HBA_i, d_{\delta})$ is a
{\bf homological model} of $A=\tilde{\ot}_{i \in I} A_i$.

We deal with the transference of the coalgebra
structure from the reduced bar construction of a CDGA
(which has a canonical coalgebra structure) to the
homological model described above.

   Let us consider the homological model $HBA$ for a CDGA
$A$. That means that we have an algebra contraction:

$$(f,g,\phi):\bar{B}(A) \stackrel{C} {\Rightarrow}
HBA$$

Let $\Delta$ be the natural coproduct on $\bar{B}(A)$:
$$\Delta([a_1|\cdots|a_n])=
\sum_{i=0}^{n}[a_1|\cdots|a_i]\ot[a_{i+1}|\cdots|a_n]$$

This coproduct induces, in a natural way, a morphism on the
homological model of $A$, that is,

\begin{equation}
\label{e1}
\Delta_2=(f\ot f)\Delta g:HBA \rightarrow HBA\ot HBA
\end{equation}

We treat the problem of determining if $\Delta_2$
is a real coproduct on the homological
model. That means that we analyze whether $\Delta_2$ is coassociative
or not.

     At first, it would be  necessary to evaluate $\Delta_2$ over all the
module generators of $HBA$. Nevertheless, since the
projection $f$ is a morphism of CDGAs (see \cite{HCDGA}), it is
easy to prove that $\Delta_2$ satisfies the following condition:

$$ \Delta_2 \circ \mu = (\mu \ot \mu)(1\ot T\ot 1) (\Delta_2 \ot \Delta_2),$$
where $\mu: HBA\ot HBA \rightarrow HBA$ is the product of the CDGA
$HBA$ and $T:HBA \ot HBA \rightarrow HBA\ot HBA$ is the morphism
that interchanges the factors. In particular, that relation means that for
determining the morphism $\Delta_2$, we
only need to evaluate it over the generators of $HBA$ as an
algebra (a finite number of elements!). Moreover, we attack the
problem of
the complexity of the associated algorithm to the explicit
formula of $\Delta_2$ and we considerably reduce the amount of
elementary operations (comparing it to the naive algorithm derived for
the initial formulation (\ref{e1})) that are necessary in order to
compute this
morphism over an algebra generator.



    Finally, we prove that if $\Delta_2$ is coassociative for the
    algebra generators of $HBA$, then $\Delta_2$ is coassociative
    for any element of $HBA$. This gives us the possibility to
    design an algorithm testing the coassociativity of this
    morphism. Several examples ``tracing''
    the previous algorithms in
    concrete cases are given.

\begin{thebibliography}{99}

\bibitem{casc98}V. \'{A}lvarez, A. Armario, R. Gonz\'{a}lez-D\'{\i}az and
P.Real. {\it Algorithms in Algebraic Topology and Homological
Algebra: the problem of the complexity}.
Workshop on Computer Algebra in Scientific Computing, CASC'98.
Saint Petersburg (April 1998).

\bibitem{Moscu} V. \'{A}lvarez, A. Armario, P. Real and B. Silva.
{\it Homological Perturbation Theory and computability of
Hochschild and cyclic homologies of CDGAs}. The international
Conference on Secondary Calculus and Cohomological Physics. EMIS
Electronic Proceedings. http://www.emis.de/proceedings/SCCP97

\bibitem{HCDGA}V. \'{A}lvarez, A. Armario, P. Real and B. Silva.
{\it Computability of the algebra homology of CDGAs}. In
preparation.


\bibitem{NY} A. Armario, P. Real and B. Silva. {\it On $p$-minimal
homological models of TTPs of elementary complexes at a prime}.
Contemporary Mathematics, {\bf vol. 227} (1999), 303-314.

\bibitem{EM53} S.  Eilenberg   and S. Mac
Lane. {\em On the groups $H(\pi,n)$, I}. Annals of Math., \mbox{\bf v.
 58} (1953), 55-106 .

\bibitem{EM54} S. Eilenberg and S. Mac Lane. {\em On the groups
$H(\pi,n)$, II}. Annals of Math., \mbox{\bf v. 60 } (1954), 49-139.



\bibitem{Gug72} V. K. A. M. Gugenheim. {\it On the cahin complex of a
fibration}. Illinois J. Math. {\bf 3} (1972) 398-414.

\bibitem{GL89} V. K. A. M. Gugenheim and L. Lambe. {\it Perturbation
theory in Differential Homological Algebra, I}. Illinois J. Math
{\bf 33} (1989) 56-582.

\bibitem{GLS91} V. K. A. M. Gugenheim, L. Lambe and J. Stasheff.
{\it Perturbation theory in Differential Homological Algebra, II}.
Illinois J. Math. {\bf 35} n. 3 (1991) 357-373.

\bibitem{HK} J. Huebschmann and T. Kadeishvili. {\em Small models for
chain algebras}. Math. Zeit. {\bf v. 207}, pp. 245-280 (1991).


\bibitem{Lam92} L. Lambe. {\em Homological perturbation theory,
Hochschild homology and formal group}.
Deformation theory and quantum groups with applications to
mathematical physics (Amherst, MA, 1990), Contemp. Math.
{\bf v. 134}, 183-218. Amer. Math. Soc., Providence, RI, 1992.


\bibitem{MCL} S. Mac Lane. {\em Homology}. Classics in
Mathematics. Springer-Verlag, Berlin (1995). Reprint of the 1975
edition.


\bibitem{Rea96a} P. Real. {\em Homological Perturbation Theory and
associativity}. Preprint del Dpto. de Mat. Apl. I.


\bibitem{RUBIO} J. Rubio. {\em Homologie effective des espaces de
lacets it\'{e}r\'{e}s: un logiciel}. PhD. thesis, Universit\'{e} Joseph
Fourier, Grenoble (1991).

\bibitem{Wei94} C. A. Weibel. {\em An introduction to homological
algebra.} Cambridge studies in advanced mathematics, {\bf 38},
Cambridge University Press (1994).





\end{thebibliography}
\end{document}

